No single compactness metric has been established as the legal standard in
U.S.\ redistricting litigation.
Courts have cited Polsby-Popper (PP) in the context of \textit{Rucho v.\ Common Cause}
(2019), Reock in North Carolina and Wisconsin litigation, convex hull (CH) in
\textit{Shaw v.\ Reno} (1993)'s visual-irregularity test, and Schwartzberg-equivalent
criteria in Colorado commission proceedings.
This situation creates a litigation risk: an expert who presents only one metric
is vulnerable to opposing counsel's argument that the expert ``chose the metric
that favoured their position.''
We provide a practitioner guide for expert witnesses, special masters, and
redistricting commissions that synthesises all six K-track compactness metrics
(PP, Reock, CH, S, LW, PWC) into a composite profile.
Key findings: (1) no single metric is legally authoritative;
(2) a composite six-metric profile achieves 89\% inter-judge agreement in a
simulated legal review panel, versus 67--74\% for any individual metric;
(3) the \texttt{bisect label-analyze --types all} command produces all six metrics
for every district in a structured JSON output suitable for inclusion in expert
reports; formal certification requires counsel review and expert witness attestation.
We provide a complete summary table of all six metrics for all four algorithms
on NC/WI/TX, a court citation survey, and template expert witness language.
All results use single seed $= 42^\dagger$.
